\chapter{Introdução}

\section{Definição do problema}

Descrever o problema central que motivou o trabalho, delimitando o contexto
em que ele ocorre e os efeitos observados.

\section{Justificativa}

Apresentar a relevância do tema e por que a investigação proposta é
necessária no cenário estudado.

\section{Objetivo geral}

Indicar, de forma direta, o que o trabalho pretende alcançar ao final da
pesquisa.

\section{Objetivos específicos}

\begin{itemize}
	\item Detalhar os passos necessários para atingir o objetivo geral;
	\item Organizar a análise do problema em partes menores;
	\item Preparar a base para a proposta apresentada nos capítulos seguintes.
\end{itemize}

\section{Estrutura do trabalho}

Apresentar, de forma resumida, como o documento está organizado:
capítulo 1 com a introdução, capítulo 2 com a fundamentação teórica,
capítulo 3 com a proposta, capítulo 4 com ferramentas e métodos e
capítulo 5 com o encerramento.
